% ##############################################################################
\section{How an intelligence market creates value}
\label{sec:market_design}

To motivate the design choices made in the rest of the paper, we step back
and consider the participants, roles, and the central design difficulty of
a market for machine intelligence.

% ==============================================================================
\subsection{Market participants}
\label{sec:market_participants}

The demand side of the market is made of applications and autonomous agents
that consume LLM inference as an input: a customer-support agent, a
document-summarization pipeline, a coding assistant, or a multi-agent
system delegating a sub-task. These participants differ along several
dimensions:

\begin{itemize}
  \item How much they value low latency.
  \item How much they value a low error rate.
  \item How price-sensitive they are.
  \item How much of their traffic is predictable in advance versus bursty.
\end{itemize}

The supply side is made of model and compute providers: foundation-model
API vendors, open-weight model hosts, and resellers of spare inference
capacity. Sellers differ along several dimensions:

\begin{itemize}
  \item The capability tier of the models they can serve.
  \item Their typical latency and reliability under load.
  \item Their marginal cost of serving a request, which itself depends on
    hardware utilization and contractual commitments elsewhere.
\end{itemize}

As in any market, the opportunity to trade arises from these differences:
a buyer with a latency-insensitive, cost-sensitive batch job and a seller
with idle capacity outside of peak hours can both be made better off by a
contract that a rigid, single posted price cannot express.

% TODO(ai_gp): Add a figure showing the two-sided market: demand-side
% dimensions (latency sensitivity, error tolerance, price sensitivity,
% traffic predictability) on one side, supply-side dimensions (capability
% tier, latency/reliability under load, marginal cost) on the other, with
% NoesisMarket as the meeting point in between.

% ==============================================================================
\subsection{Roles of an intelligence market}
\label{sec:market_roles}

An intelligence market, like any exchange, has to perform several
functions beyond simply connecting a buyer to a seller:

\begin{itemize}
  \item Provide a common meeting ground where buyers and sellers can
    express demand and supply without prior bilateral negotiation.
  \item Establish a standardized contract schema (Section~\ref{sec:contract})
    so that a bid and an ask submitted by unrelated parties are directly
    comparable.
  \item Provide oversight by metering whether a matched contract was
    actually fulfilled, a role played in \Noesis{} by \NoesisServer{}
    rather than by a separate compliance function.
  \item Generate reusable market data: clearing prices per tier and round,
    and, as a byproduct of fulfillment monitoring, a record of measured
    provider performance over time.
\end{itemize}

As in traditional exchanges, we expect the market to be compensated for
this value creation through a small fee on cleared volume; we do not
formalize a fee schedule in this paper.

% ==============================================================================
\subsection{The bundling problem: quoting quality, not just quantity}
\label{sec:bundling}

The central design difficulty is choosing what, precisely, is being
bought and sold. A market for a homogeneous good only needs to clear on
price and quantity. Machine intelligence is not homogeneous: the same
nominal price per request can correspond to wildly different quality, and
the same underlying compute can be used to answer the same request at very
different quality depending on the model and routing policy used to serve
it.

This is analogous to the base/quote asymmetry that a token-exchange
protocol must resolve before a limit order is even well-defined: a limit
order only makes sense once one has fixed which token plays the role of
quantity and which plays the role of price. \Noesis{} faces a similar
prior question: before a bid or an ask can be written down, the protocol
must fix what a unit of ``intelligence'' is (Section~\ref{sec:task_unit})
and what accompanies that unit as a quality specification. We adopt the
position that quality should be expressed along three largely orthogonal
axes:

\begin{itemize}
  \item \textbf{Capability}: how good the answer is expected to be,
    expressed as a discrete tier rather than a continuous score, since
    tiers are what providers can credibly commit to and what buyers can
    reason about.
  \item \textbf{Latency}: an upper bound on the time to produce an answer.
  \item \textbf{Reliability}: the fraction of requests that meet both the
    capability and latency terms.
\end{itemize}

% TODO(ai_gp): Add a figure showing the capability/latency/reliability bundle
% as three axes of a single traded good, contrasted with a single-scalar
% price-per-token good, to make the bundling problem visually concrete.

Section~\ref{sec:contract} turns this into a formal contract. The
consequence of this design choice is that the market cannot clear on price
alone: two asks at the same price but different tiers are not
substitutes, and matching must first establish compatibility along these
three axes before price discovery can proceed
(Section~\ref{sec:compatibility}).

% ==============================================================================
\subsection{Related commodity markets}
\label{sec:related_commodity_markets}

Electricity markets are the closest existing analogy. Electricity cannot
be usefully stored at scale, demand and supply must be balanced
continuously, and quality (voltage, frequency) is metered rather than
taken on faith. Electricity markets address this with a layered design:
day-ahead markets clear the bulk of expected demand a day in advance,
real-time markets clear residual imbalances close to delivery, and
separate capacity markets pay generators for being available rather than
for energy delivered, penalizing non-delivery when it matters most
\cite{stoft2002, cramton2017}. \Noesis{}'s tiered call auction
(Section~\ref{sec:noesis_market}) borrows the day-ahead/real-time
separation as a candidate design (discussed further as an open question in
Section~\ref{sec:open_questions}), and its reputation and pricing feedback
loop (Section~\ref{sec:reputation}) borrows directly from capacity-market
non-performance penalties \cite{cramtonstoft2005}.

The analogy is imperfect in one important respect: electricity's quality
attributes (voltage, frequency) are physical and essentially free to
meter, while machine intelligence's capability attribute is not directly
observable and can only be estimated from the content of a response. This
is the problem that \NoesisServer{}'s fulfillment monitoring
(Section~\ref{sec:fulfillment_monitoring}) has to solve.

% TODO(gp): Add a timeline figure comparing electricity market layering
% (day-ahead / real-time / capacity markets) against Noesis's periodic
% tiered call auction, to make the borrowed structure concrete.
